
\subsection{Różniaste różności}
\paragraph{Warunki Torresa}
Istnieje odmiana wielomianu Alexandera, która liczy sobie tyle zmiennych, ile ogniw posiada splot (nie opisywaliśmy jej i~nie zamierzamy).
Klasyfikację \ref{prp:alexander_hosokawa} można częściowo uogólnić: Torres \cite{torres1953} znalazł dwie geometryczne własności, nazwane później warunkami Torresa.
\index[persons]{Torres, Guillermo}%
\index{warunek!Torresa}%
Są warunkami koniecznymi, ale nie wystarczającymi, jak odkrył ponad ćwierć wieku później Hillman \cite{hillman1981}: wielomian
\index[persons]{Hillman, Jonathan}%
\begin{equation}
    D(x,y) = \frac{(1 - x^6y^6)(x - 1 + 1/x) - 2(1 - x^5y^5)(1 - x)(1 - y)}{1-xy}
\end{equation}
spełnia warunki Torresa, ale nie jest wielomianem Alexandera żadnego splotu.
Wspominają o tym krótko Burde, Zieschang \cite[s. 150]{burde2014}.

\paragraph{Unimodalne współczynniki}
Fox \cite{fox1962} podejrzewał, że
\index[persons]{Fox, Ralph}%
\index{hipoteza!trapezoidalna}%
ciąg współczynników wielomianu Alexandera węzła alternującego jest unimodalny.
Dowód podano dla węzłów algebraicznych (Murasugi \cite{murasugi1985}) oraz genusu dwa (Ozsváth i~Szabó \cite{ozsvath2003}).
\index[persons]{Murasugi, Kunio}%
\index[persons]{Ozsváth, Peter}%
\index[persons]{Szabó, Zoltán}%
Hipoteza w~ogólnym przypadku pozostaje otwarta.
Sama hipoteza Foxa \cite[problem 12]{fox1962} jest nieco mocniejsza.
Niech $\alexander_K(t) = \sum_{i=0}^n a_i t^i$ (z~dokładnością do jednomianu) będzie wielomianem Alexandera alternującego węzła $K$, wtedy $|a_i| = |a_{n-i}|$.
Hipoteza trapezoidalna głosi, że
\begin{equation}
    |a_0| \le |a_1| \le \ldots \le |a_{\lfloor n/2 \rfloor}|,
\end{equation}
oraz że jeśli $|a_i| = |a_{i+1}|$ dla pewnego $i < n/2$, to $|a_i| = |a_{i+1}| = \ldots = |a_{n-i}|$.
Wykres współczynników ma więc kształt trapezu z~płaskim wierzchołkiem.
Stojmenow \cite{stoimenowb2005} podejrzewa więcej: dla alternujących splotów ciąg $|a_i|$ jest logarytmicznie wklęsły, to znaczy $|a_{i-1}| \, |a_{i+1}| \le a_i^2$ dla $0 < i < n$; dla $i < n/2$ równość zachodzi tylko wtedy, gdy $a_{i-1} = -a_i = a_{i+1}$, a~długość płaskiego wierzchołka nie przekracza wartości bezwzględnej sygnatury.
\index{hipoteza!trapezoidalna}%
\index{wielomian!log-wklęsły}%
Ta mocniejsza wersja jest znana dla węzłów dwumostowych (Banfield \cite{banfield2022}) oraz dla splotów specjalnie alternujących, to znaczy mających alternujący diagram, w~którym wszystkie skrzyżowania mają ten sam znak (Hafner, Mészáros, Vidinas \cite{hafner2024}).
Stojmenow sprawdził ją dla alternujących węzłów genusu co najwyżej $4$ \cite[sekcja 9.2]{stoimenow2016}.
\index[persons]{Banfield, Ian}%
\index[persons]{Hafner, Elena}%
\index[persons]{Mészáros, Karola}%
\index[persons]{Vidinas, Alexander}%
Dalsze postępy w~stronę hipotezy Foxa opisują Azarpendar, Juhász, Kálmán \cite{azarpendar2024}; ogólne omówienie znajduje się w~\cite[problem 1.7]{baykur2026}.
\index[persons]{Azarpendar, Soheil}%
\index[persons]{Juhász, András}%
\index[persons]{Kálmán, Tamás}%

Dla dodatnich splotów o~$\mu$ ogniwach wielomian Conwaya ma postać $\conway_L(z) = z^{\mu-1} \sum_{i=0}^m c_i z^{2i}$ z~dodatnimi współczynnikami $c_i$.
Stojmenow przypuszcza, że ciąg $(c_i)$ jest silnie logarytmicznie wklęsły, to znaczy $c_{i-1} c_{i+1} < c_i^2$ dla $0 < i < m$; dla splotów specjalnie alternujących udowodnił to w~\cite[tw. 1]{stoimenowb2005}.
Wiemy natomiast, że kolejne współczynniki wielomianu Conwaya węzła alternującego są przeciwnych znaków i niezerowe (Murasugi \cite[s. 242]{murasugi1996} odsyła do swojej wcześniejszej pracy \cite{murasugi1959}, gdzie używa archaicznej nomenklatury: \textsc{Schlauchknoten} zamiast węzłów satelitarnych jak w \cite[s. 245]{schubert1953}).
\index{Schlauchknoten}%
Wynika stąd, że węzły $(p, q)$-torusowe dla $p > q > 2$ oraz pierwsze węzły satelitarne nie są alternujące.

% koniec podsekcji Różniaste różności

